\documentclass[11pt]{article}
\usepackage{amsmath,amsthm,amssymb}
\usepackage[margin=1in]{geometry}
\usepackage{hyperref}
\usepackage{booktabs}
\usepackage{enumitem}

\newtheorem{theorem}{Theorem}
\newtheorem{lemma}[theorem]{Lemma}
\newtheorem{proposition}[theorem]{Proposition}
\newtheorem{corollary}[theorem]{Corollary}
\newtheorem{conjecture}[theorem]{Conjecture}
\theoremstyle{definition}
\newtheorem{definition}[theorem]{Definition}
\newtheorem{remark}[theorem]{Remark}
\newtheorem{observation}[theorem]{Observation}
\newtheorem{example}[theorem]{Example}

\newcommand{\Hq}{H_q}
\newcommand{\Sn}{S_n}
\newcommand{\Vlam}{V_\lambda}
\newcommand{\Bdec}{B_{\ell+1}^{(\lambda)}}
\newcommand{\Om}{\Omega}
\newcommand{\hatl}{\widehat\lambda}
\newcommand{\supp}{\mathrm{supp}}
\newcommand{\SYT}{\mathrm{SYT}}
\newcommand{\Rp}{{R'}}

\title{Extended sign-kill for $\hatl=(4,4)$: \\
       single Phi-branch with same-row $(4,2)$ recursion}
\author{Clio}
\date{14 May 2026 (very late prove session, ninth pass)}

\begin{document}
\maketitle

\begin{abstract}
Let $\lambda = (4, 4, 1^q) \vdash n = q+8$ with $q \ge 6$,
$\ell = q+2 = n-6$, and $B = \Bdec\Vlam = \Rp_{\ell+1}\Rp_{\ell+2}\cdots\Rp_n\,\Vlam$
($\Om$ has \emph{six} $R'$-blocks).  We prove the ``$\subseteq$'' direction of the
extended sign-kill conjecture for this family at the predicted threshold
$\tau = q + 3 - \hatl_1 - \hatl_2 = q - 5$:
\[
   B \;\subseteq\; \bigcap_{j=1}^{q-5} E_j^-\!\mid_{\Vlam}.
\]
At $q \le 5$ the predicted threshold is non-positive, so the containment
is vacuous.

The shape $\hatl = (4, 4)$ is the second $r=1$ multi-corner case
in the programme (after $\hatl = (3, 3)$).  Only two corners
$c_1 = (2, 4)$ and $c_2 = (q+2, 1)$ are removable, so the
$E_{n-1}^+$ decomposition splits into a single $2$-block piece
$\Phi_A(V_{\mu_A})$ over $\mu_A = (4, 3, 1^{q-1})$ plus a same-row
piece $\Phi_h(V_{\mu_h})$ over $\mu_h = (4, 2, 1^q)$.  Both inner
shapes are non-hooks, so $\hatl = (4, 4)$ is the first $r=1$ case
where the same-row vanishing requires a (non-hook) extended sign-kill
input --- namely the May-14 sixth-pass $(4, 2)$ theorem.  The
contributing branch recurses on the May-14 eighth-pass $(4, 3)$ theorem.

As a corollary, the proof extracts a structural ``stability ladder''
for the $(4, 4, 1^q)$ family:
\begin{itemize}[topsep=0.2em,itemsep=0.1em]
\item $q \in \{2, 3\}$: \emph{sub-stable}.  Both $\Phi_A$ and $\Phi_h$
contribute; $\dim B = 10, 15$ (factor $2, 3$ of the stable value $5$).
\item $q \in \{4, 5\}$: \emph{stable but pre-threshold}.  $\Phi_h$
vanishes, $\dim B = 5$, but the recursive $(4, 3)$ input is not yet
available, so the prefix bound and $E_j^-$ chain are weaker than the
predicted $\tau$.
\item $q \ge 6$: \emph{full regime}.  $\dim B = 5$, $B \subseteq E_j^-$
for $1 \le j \le q - 5$, with sharpness at $j = q - 4$.
\end{itemize}
\end{abstract}

\section{Setup}\label{sec:setup}

We use the conventions of the May-13--May-14 series.  $\Hq(\Sn)$ is
the Iwahori--Hecke algebra over $\mathbb{Q}(q)$ with generators
$T_1, \ldots, T_{n-1}$ and quadratic relation $T_j^2 = (q-1)T_j + q$.
$\Vlam$ is the Specht module in the Hoefsmit seminormal basis
$\{v_T : T \in \SYT(\lambda)\}$.

For $1 \le j \le n-1$, set $S_j := T_j + 1$.  Write
\[
   E_j^+ \;:=\; \ker(T_j - q)\!\mid_{\Vlam}
   \;=\; \mathrm{Im}(S_j\!\mid_{\Vlam}),
   \qquad
   E_j^- \;:=\; \ker(S_j)\!\mid_{\Vlam}.
\]
For $k \ge 1$, let $\Rp_k := S_{k-1}\,S_{k-2}\cdots S_1$ with the
convention $\Rp_1 = 1$.  For $\lambda \vdash n$ with
$\ell = \ell(\lambda)$, set
\[
   \Om \;=\; \Om^{(\lambda)} \;:=\; \Rp_{\ell+1}\Rp_{\ell+2}\cdots \Rp_n,
   \qquad
   B \;:=\; \Om \cdot \Vlam \;=\; \Bdec \Vlam.
\]

For $\lambda = (4, 4, 1^q)$, $n = q + 8$ and $\ell = q + 2$, so
\[
   \Om \;=\; \Rp_{q+3}\Rp_{q+4}\Rp_{q+5}\Rp_{q+6}\Rp_{q+7}\Rp_{q+8}
   \;=\; \Rp_{n-5}\Rp_{n-4}\Rp_{n-3}\Rp_{n-2}\Rp_{n-1}\Rp_n
\]
has \emph{six} $\Rp$-blocks.  By the May-13 multi-corner-dimension
paper, $\dim B = f^{\lambda^{(\ge 2)}} = f^{(3, 3)} = 5$ in the stable
regime.  For this family, computation (mod $P = 100\,003$, $Q = 11$)
shows that the stable regime begins at $q \ge 4$: at $q = 2$ we get
$\dim B = 10 = 2{\cdot}5$, at $q = 3$ we get $\dim B = 15 = 3{\cdot}5$
(both \emph{sub-stable}, related to the row-repeat factor effect noted
in the May-14 evening factor-of-two refutation paper), and at $q \ge
4$ we observe $\dim B = 5$ stably.  Our main theorem (Section~\ref{sec:main})
assumes $q \ge 6$ for a stronger reason --- the recursive $(4, 3)$
input requires $q' = q - 1 \ge 5$.

\paragraph{Removable corners of $\lambda$.}
Since $\lambda_1 = \lambda_2 = 4$, the cell $(1, 4)$ is \emph{not}
removable.  The removable corners are
\[
   c_1 = (2, 4), \qquad c_2 = (q+2, 1).
\]
There are only \emph{two} corners, so $\lambda$ is an $r = 1$ shape in
the multi-corner classification.

\paragraph{Identity prefix.}  For an SYT $T$ of any partition $\mu$,
the \emph{identity prefix length} is
\[
   p(T) \;:=\; \max\{\,P \ge 0 : T(i, 1) = i \text{ for all }
   1 \le i \le P\,\}.
\]

\paragraph{Seminormal support.}  For $W \subseteq V_\lambda$,
\[
   \supp(W) \;:=\; \{\,T \in \SYT(\lambda) : \text{some } w \in W
   \text{ has } [v_T]\,w \ne 0\,\}.
\]

\section{Input from prior papers}\label{sec:input}

\subsection{Support rule and prefix corollary (May-14 fourth pass)}

\begin{proposition}[Support rule]\label{prop:support-rule}
Let $W \subseteq V_\lambda$ be any subspace.  Fix $1 \le j \le n-1$.
If every $T \in \supp(W)$ has $j$ and $j+1$ in the same column of $T$,
then $W \subseteq E_j^-\!\mid_{V_\lambda}$.
\end{proposition}

\begin{remark}\label{rem:prefix-implies-Ej-}
If every $T \in \supp(W)$ has $p(T) \ge P$, then for every
$1 \le j \le P - 1$ the entries $j, j+1$ sit at column-$1$ rows
$j, j+1$, sharing column~$1$.  By Proposition~\ref{prop:support-rule},
$W \subseteq E_j^-$ for $1 \le j \le P - 1$.
\end{remark}

\subsection{Outer-factor preservation (May-14 fifth pass)}

\begin{lemma}[Outer-factor preservation]\label{lem:outer-preserve}
Let $W \subseteq V_\nu$ be a subspace such that every $T \in \supp(W)$
satisfies $p(T) \ge P$.  Then for every $j \ge P + 1$,
every $T'' \in \supp((T_j{+}1) W)$ also satisfies $p(T'') \ge P$.
\end{lemma}

\begin{proof}[Proof sketch]
$(T_j+1)$ acts within the $2$-block $\{v_T, v_{s_j T}\}$; the swap $s_j$
moves only entries $j$ and $j+1$, both $\ge j \ge P + 1 > P$, so neither
lies in the column-$1$ prefix and the prefix is preserved.\qed
\end{proof}

\subsection{The May-14 eighth-pass $(4, 3)$ extended sign-kill theorem}

\begin{theorem}[Extended sign-kill for $\hatl = (4, 3)$, eighth pass]\label{thm:43-prefix}
For $\lambda' = (4, 3, 1^{q'}) \vdash n' = q' + 7$ with $q' \ge 5$,
every $T' \in \supp\bigl(B^{(\lambda')}_{q'+3} V_{\lambda'}\bigr)$
satisfies $p(T') \ge q' - 3 = n' - 10$.  Consequently,
$B^{(\lambda')}_{q'+3} V_{\lambda'} \subseteq E_j^-\!\mid_{V_{\lambda'}}$
for $1 \le j \le q' - 4$, and in particular $B \subseteq E_1^-$.
\end{theorem}

\subsection{The May-14 sixth-pass $(4, 2)$ extended sign-kill theorem}

\begin{theorem}[Extended sign-kill for $\hatl = (4, 2)$, sixth pass]\label{thm:42-prefix}
For $\lambda' = (4, 2, 1^{q'}) \vdash n' = q' + 6$ with $q' \ge 4$,
every $T' \in \supp\bigl(B^{(\lambda')}_{q'+3} V_{\lambda'}\bigr)$
satisfies $p(T') \ge n' - 8 = q' - 2$.  Consequently,
$B^{(\lambda')}_{q'+3} V_{\lambda'} \subseteq E_j^-\!\mid_{V_{\lambda'}}$
for $1 \le j \le q' - 3$, and in particular $B \subseteq E_1^-$.
\end{theorem}

\subsection{Phase A endpoint (May-19)}

\begin{theorem}[Phase A endpoint]\label{thm:phaseA}
For any $\Hq$-module $V$ and any $1 \le j \le n-1$,
$(T_j+1)\cdots(T_1+1) V = E_j^+(V)$.  In particular,
$\Rp_n V_\lambda = E_{n-1}^+|_{V_\lambda}$.
\end{theorem}

\section{The $+q$-eigenspace decomposition for $(4,4,1^q)$}\label{sec:Eplus}

We give the structure of $\Rp_n V_\lambda = E_{n-1}^+|_{V_\lambda}$
analogously to the $\hatl = (3, 3)$ paper.

\paragraph{2-block piece (doubly-stripped at $\{c_1, c_2\}$).}
The pair $\{c_1, c_2\} = \{(2, 4), (q+2, 1)\}$ has axial distance
\[
   |d| \;=\; |c(c_1) - c(c_2)|
   \;=\; |2 - (-q-1)| \;=\; q + 3 \;\ge\; 5
   \quad \text{for } q \ge 2,
\]
so the corresponding $T_{n-1}$-block is non-degenerate.  The
doubly-stripped partition is
\[
   \mu_A \;:=\; \lambda \setminus \{c_1, c_2\} \;=\; (4, 3, 1^{q-1})
   \quad (\text{a $\hatl = (4, 3)$ shape}).
\]

\paragraph{Same-row piece at $c_1$.}
The cell $(2, 3)$ sits immediately left of $c_1 = (2, 4)$ in row 2.
SYTs of $\lambda$ with $n-1$ at $(2, 3)$ and $n$ at $(2, 4)$ are
$+q$-eigenvectors of $T_{n-1}$ (same-row), distinct from the
$2$-block pieces.  Their inner shape is
\[
   \mu_h \;:=\; \lambda \setminus \{(2, 3), (2, 4)\}
   \;=\; (4, 2, 1^q)
   \quad (\text{a $\hatl = (4, 2)$ shape}).
\]
Note that $\mu_h$ is \emph{not} a hook (in contrast to the $(3, 3)$
case): row 2 of $\mu_h$ still has length 2.  Hence the same-row
vanishing in the present paper does \emph{not} reduce to the May-12
hook column-$1$ theorem; it requires the $(4, 2)$ extended sign-kill
input (Theorem~\ref{thm:42-prefix}).

\paragraph{No other +q pieces.}
Same-row at $c_2 = (q+2, 1)$: no cell to its left, so vacuous.
Same-column extensions at column $1$ produce $-1$-eigenvectors
($E_{n-1}^-$), not relevant here.  Cells $(1, k)$ in row 1 with
$k \le 4$ cannot be terminal positions for $n$ because the cell
below them, $(2, k)$, still exists in $\lambda$.

\begin{lemma}[Two-piece decomposition of $E_{n-1}^+$]\label{lem:Eplus-decomp}
Define $\Phi_A : V_{\mu_A} \hookrightarrow V_\lambda$ on the
seminormal basis by
\[
   \Phi_A(v_{T'}) \;:=\; v_{T_A^A} + \tau_A\,v_{T_A^B},
\]
where $T_A^A, T_A^B$ are the two SYTs of $\lambda$ obtained from
$T' \in \SYT(\mu_A)$ by placing $\{n-1, n\}$ at $\{c_1, c_2\}$ in the
two possible orders, and $\tau_A \in \mathbb{Q}(q)^\times$ is the
Hoefsmit 2-block coefficient at axial distance $|d| = q+3$.

Define $\Phi_h : V_{\mu_h} \hookrightarrow V_\lambda$ on the seminormal
basis by $\Phi_h(v_{T''}) := v_{T'''}$, where $T'''$ is the unique
SYT of $\lambda$ obtained from $T'' \in \SYT(\mu_h)$ by placing $n-1$
at $(2, 3)$ and $n$ at $(2, 4)$.

Then:
\begin{enumerate}[topsep=0.3em,itemsep=0.1em,leftmargin=2em,
                  label=(\roman*)]
\item Both $\Phi_A$ and $\Phi_h$ are injective and $T_i$-equivariant
   for $1 \le i \le n-3$.
\item Their images in $V_\lambda$ have disjoint seminormal supports
   and satisfy
   \[
       \Rp_n V_\lambda \;=\; E_{n-1}^+\!\mid_{V_\lambda}
       \;=\; \Phi_A(V_{\mu_A}) \;\oplus\; \Phi_h(V_{\mu_h}).
   \]
\end{enumerate}
\end{lemma}

\begin{proof}
Injectivity and $T_i$-equivariance for $1 \le i \le n-3$ are
standard: for $i \le n-3$ the swap $s_i$ moves only entries $i, i+1
\le n - 2$, which live inside the ``inner'' part ($\mu_A$- or
$\mu_h$-cells); the Hoefsmit coefficients depend only on those
inner cells, so $T_i$ commutes with the corresponding $\Phi$.

Disjointness of supports: $\Phi_A$ supports require $\{n-1, n\}$ at
$\{(2, 4), (q+2, 1)\}$ (one entry each); $\Phi_h$ supports require
$n-1 = (2, 3)$ and $n = (2, 4)$.  These cell configurations are
distinct (position of $n-1$ differs), so supports are disjoint.

Dimensions: $\dim \Phi_A(V_{\mu_A}) = f^{(4, 3, 1^{q-1})}$ and
$\dim \Phi_h(V_{\mu_h}) = f^{(4, 2, 1^q)}$.  Enumerating $+q$-eigenvectors
of $T_{n-1}$ on $V_\lambda$:
\[
   \dim E_{n-1}^+ \;=\; \#\bigl(\text{2-block pairs at }\{c_1, c_2\}\bigr)
   + \#\bigl(\text{same-row SYTs at }c_1\bigr)
   \;=\; f^{(4, 3, 1^{q-1})} + f^{(4, 2, 1^q)}.
\]
The two pieces are disjoint subspaces of $E_{n-1}^+$ with summed
dimension equal to $\dim E_{n-1}^+$, hence span (in direct sum) all
of $E_{n-1}^+$.  By Theorem~\ref{thm:phaseA},
$\Rp_n V_\lambda = E_{n-1}^+$, giving the claimed equality.

\emph{Computational verification.}  At $q \in \{1, 2, 3, 4\}$ the
dimension match
$\dim \Rp_n V_\lambda = f^{(4,3,1^{q-1})} + f^{(4,2,1^q)}$
was verified directly
(script \texttt{\textasciitilde/projects/scratch/2026-05-14-extended-sign-kill-44/check\_decomp.py}).
\qed
\end{proof}

\section{The single-branch reduction for $(4, 4, 1^q)$}\label{sec:reduction}

\begin{theorem}[Single-branch reduction]\label{thm:reduction}
For $\lambda = (4, 4, 1^q)$ with $q \ge 4$,
\begin{equation}\label{eq:reduction}
   B \;=\; (T_{n-6}{+}1)(T_{n-5}{+}1)(T_{n-4}{+}1)(T_{n-3}{+}1)(T_{n-2}{+}1)\,
   \Phi_A\bigl(B^{(\mu_A)}_{\ell(\mu_A)+1} V_{\mu_A}\bigr).
\end{equation}
At $q \in \{2, 3\}$ the same-row branch $\Phi_h$ does \emph{not}
vanish (the $(4, 2)$ extended sign-kill input requires $q' \ge 4$),
so $B$ contains both contributions; this explains the sub-stable
factor of $2$ and $3$ in $\dim B$ observed at $q = 2$ and $q = 3$.
\end{theorem}

\begin{proof}
Following the May-15 / May-20 pull-through technique, apply it
piecewise to each summand in
$\Rp_n V_\lambda = \Phi_A(V_{\mu_A}) \oplus \Phi_h(V_{\mu_h})$
(Lemma~\ref{lem:Eplus-decomp}), then show that the same-row branch
$\Phi_h$ contributes $0$.

\medskip\textbf{Step 1 (pulling the chain through the $\Phi$ pieces).}
Let $\Phi \in \{\Phi_A, \Phi_h\}$ and $\mu \in \{\mu_A, \mu_h\}$
correspondingly.

Decompose $\Rp_{n-1} = (T_{n-2}+1)\,U_1$ where
$U_1 = (T_{n-3}+1)\cdots(T_1+1)$.  Since $T_1, \ldots, T_{n-3}$ commute
with $T_{n-1}$ (index gaps $\ge 2$), $U_1$ preserves
$E_{n-1}^+|_{V_\lambda}$.  By Lemma~\ref{lem:Eplus-decomp}(i),
$U_1 \Phi(v) = \Phi(\Rp_{n-2}^{(\mu)}\,v)$ for $v \in V_\mu$.  Hence
\[
   \Rp_{n-1}\,\Phi(V_\mu) \;=\; (T_{n-2}+1)\,\Phi\bigl(\Rp_{n-2}^{(\mu)} V_\mu\bigr).
\]

Iterating: decompose $\Rp_{n-k} = (T_{n-k-1}+1)\,U_k$ where
$U_k = (T_{n-k-2}+1)\cdots(T_1+1)$.  For each $k = 2, 3, 4, 5$, the
factor $U_k$ commutes with $T_{n-2}, T_{n-3}, \ldots, T_{n-k}$
(index gaps $\ge 2$) and acts on $\Phi(v)$ via $\Phi$.  The pull-through
thus accumulates outer $(T_{n-k-1}+1)$ factors and pushes the
$\Rp^{(\mu)}_{n-k}$ factors into the inner $\mu$-action:
\begin{equation}\label{eq:branchpull}
   \Rp_{n-5}\Rp_{n-4}\Rp_{n-3}\Rp_{n-2}\Rp_{n-1}\,\Phi(V_\mu)
   \;=\;
   \prod_{k=2}^{6} (T_{n-k}+1)\;\cdot\;
   \Phi\!\Bigl(\Rp^{(\mu)}_{n-6} \Rp^{(\mu)}_{n-5} \Rp^{(\mu)}_{n-4}
                \Rp^{(\mu)}_{n-3} \Rp^{(\mu)}_{n-2}\,V_\mu\Bigr),
\end{equation}
where the outer factor product $\prod_{k=2}^{6}(T_{n-k}+1) =
(T_{n-2}+1)(T_{n-3}+1)(T_{n-4}+1)(T_{n-5}+1)(T_{n-6}+1)$ stacks left to right.

\medskip\textbf{Step 2 (combining the two pieces).}
By Lemma~\ref{lem:Eplus-decomp}(ii) and Theorem~\ref{thm:phaseA},
\begin{align*}
   B \;=\; \Om V_\lambda
   &\;=\; \Rp_{n-5}\Rp_{n-4}\Rp_{n-3}\Rp_{n-2}\Rp_{n-1}\,\Rp_n V_\lambda \\
   &\;=\; \Rp_{n-5}\Rp_{n-4}\Rp_{n-3}\Rp_{n-2}\Rp_{n-1}\,\bigl(\Phi_A(V_{\mu_A}) \oplus \Phi_h(V_{\mu_h})\bigr) \\
   &\;=\; \prod_{k=2}^6 (T_{n-k}{+}1)\,
        \bigl[\Phi_A\bigl(I_A\bigr) \;+\; \Phi_h\bigl(I_h\bigr)\bigr],
\end{align*}
where $I_X := \Rp^{(\mu_X)}_{n-6} \Rp^{(\mu_X)}_{n-5}\Rp^{(\mu_X)}_{n-4}
\Rp^{(\mu_X)}_{n-3}\Rp^{(\mu_X)}_{n-2}\,V_{\mu_X}$ (\emph{five}
$\Rp$-factors with $\mu_X$-indexing).

\medskip\textbf{Step 3 (the same-row branch $\Phi_h$ vanishes).}
For $\mu = \mu_h = (4, 2, 1^q)$, $\ell(\mu_h) = q + 2 = n - 6$,
so $\ell(\mu_h) + 1 = n - 5$.  The inner chain $I_h$ has indices
$n-6, n-5, n-4, n-3, n-2$ and so includes one $\Rp^{(\mu_h)}$-factor
\emph{below} the standard $B^{(\mu_h)}_{\ell(\mu_h)+1}$ object:
\[
   I_h \;=\; \Rp^{(\mu_h)}_{n-6}\,\bigl(\Rp^{(\mu_h)}_{n-5}\,
   \Rp^{(\mu_h)}_{n-4}\,\Rp^{(\mu_h)}_{n-3}\,\Rp^{(\mu_h)}_{n-2}\,V_{\mu_h}\bigr)
   \;=\; \Rp^{(\mu_h)}_{n-6}\,B^{(\mu_h)}_{\ell(\mu_h)+1}\,V_{\mu_h}.
\]
By Theorem~\ref{thm:42-prefix} applied to $\mu_h$ at $q' = q$
(hypothesis $q' \ge 4$, i.e., $q \ge 4$),
$B^{(\mu_h)}_{\ell(\mu_h)+1}\,V_{\mu_h} \subseteq E_1^-|_{V_{\mu_h}}$.
But the rightmost factor of $\Rp^{(\mu_h)}_{n-6}$ is $(T_1+1) = S_1$,
which annihilates $E_1^-|_{V_{\mu_h}}$.  Hence $I_h = 0$ and the
$\Phi_h$ contribution vanishes.

\medskip\textbf{Step 4 (the $\Phi_A$ piece is $B^{(\mu_A)}_{\ell(\mu_A)+1}\,V_{\mu_A}$).}
For $\mu = \mu_A = (4, 3, 1^{q-1})$, $\ell(\mu_A) = q + 1 = n - 7$,
so $\ell(\mu_A) + 1 = n - 6$.  The five-factor product
$I_A = \Rp^{(\mu_A)}_{n-6} \Rp^{(\mu_A)}_{n-5}\Rp^{(\mu_A)}_{n-4}
\Rp^{(\mu_A)}_{n-3}\Rp^{(\mu_A)}_{n-2}\,V_{\mu_A}$ has indices
running from $\ell(\mu_A)+1 = n-6$ to $n-2 = |\mu_A|$, so it is
exactly $B^{(\mu_A)}_{\ell(\mu_A)+1}\,V_{\mu_A}$.

Combining Steps 2--4 gives~\eqref{eq:reduction}.
\qed
\end{proof}

\begin{remark}\label{rem:vanishing-mechanism}
The vanishing of $\Phi_h$ is the direct analogue of the hook-piece
vanishing in the $\hatl = (3, 3)$ paper.  The only structural change
is that $\mu_h = (4, 2, 1^q)$ is \emph{not} a hook --- it has a row
of length 2 below the first row.  Consequently the inner sign-kill
input is no longer the May-12 hook column-$1$ theorem but the
extended sign-kill theorem at $\hatl = (4, 2)$ (Theorem~\ref{thm:42-prefix}).
The same ``extra rightmost $\Rp^{(\mu)}_{n-6}$ factor kills $E_1^-$'' mechanism
applies: this is the same shape-uniform branch-vanishing principle observed
since the (4, 2) and (3, 3) papers.
\end{remark}

\section{Main theorem}\label{sec:main}

\begin{theorem}[Main]\label{thm:main}
For $\lambda = (4, 4, 1^q) \vdash n = q + 8$ with $q \ge 6$,
\[
   B \;\subseteq\; \bigcap_{j=1}^{q-5} E_j^-\!\mid_{V_\lambda}.
\]
\end{theorem}

\begin{proof}
By Remark~\ref{rem:prefix-implies-Ej-}, it suffices to show that
every $T \in \supp(B)$ satisfies $p(T) \ge q - 4 = n - 12$.  We use
the reduction~\eqref{eq:reduction} together with a prefix bound on
$B^{(\mu_A)}_{\ell(\mu_A)+1}\,V_{\mu_A}$ and
Lemma~\ref{lem:outer-preserve} applied to the five outer factors.

\medskip\textbf{Step 1: prefix bound on the $\Phi_A$-image.}
$\mu_A = (4, 3, 1^{q-1})$ has size $n - 2$, $\ell(\mu_A) = q + 1
= n - 7$, and trailing-$1$-count $q' = q - 1$.  Since $q \ge 6$ we
have $q' \ge 5$, so Theorem~\ref{thm:43-prefix} applies:
\[
   p(T') \;\ge\; q' - 3 \;=\; q - 4
   \qquad \text{for every } T' \in \supp(B^{(\mu_A)}_{\ell(\mu_A)+1}\,V_{\mu_A}).
\]

Each $T' \in \supp(B^{(\mu_A)}_{\ell(\mu_A)+1}\,V_{\mu_A})$ contributes
to $\Phi_A(\cdot)$ as $v_{T_A^A} + \tau_A\,v_{T_A^B}$, where $T_A^A,
T_A^B$ are the two SYTs of $\lambda$ obtained from $T'$ by placing
$\{n-1, n\}$ at $\{c_1, c_2\} = \{(2, 4), (q+2, 1)\}$ in the two
orders.  Column-$1$ cells of $\mu_A$ occupy rows $1, \ldots,
\ell(\mu_A) = n - 7$.  The added cell $(2, 4)$ is in column $4$
(not column $1$), and the added cell $(q+2, 1)$ occupies a new
column-$1$ row $q + 2 = n - 6 > n - 7 = \ell(\mu_A)$, \emph{below}
all existing column-$1$ cells of $\mu_A$.  Hence the column-$1$
identity-prefix of $T'$ is preserved exactly:
\[
   p(T_A^A),\, p(T_A^B) \;\ge\; p(T') \;\ge\; q - 4 \;=\; n - 12.
\]

\medskip\textbf{Step 2: outer factors preserve the prefix.}
The five outer factors in~\eqref{eq:reduction} are $(T_j+1)$ for
$j \in \{n-6, n-5, n-4, n-3, n-2\}$.  Each index satisfies $j \ge
n - 6 = (n - 12) + 6 > n - 11 = (n - 12) + 1$ (more weakly,
$j \ge n - 6 \ge P + 1$ with $P = n - 12$), so the hypothesis of
Lemma~\ref{lem:outer-preserve} is met for each factor.  Applying the
lemma five times, the prefix bound $\ge n - 12 = q - 4$ is preserved.

\medskip\textbf{Conclusion.}  By Steps 1--2 and~\eqref{eq:reduction},
every $T \in \supp(B)$ has $p(T) \ge q - 4$.  By
Remark~\ref{rem:prefix-implies-Ej-}, $B \subseteq E_j^-$ for
$1 \le j \le (q - 4) - 1 = q - 5$.
\qed
\end{proof}

\begin{remark}[$q \in \{2, 3, 4, 5\}$]\label{rem:q-low}
At $q \le 5$ the predicted threshold $\tau = q - 5 \le 0$, so the
intersection $\bigcap_{j=1}^{\tau} E_j^-$ is the whole module and the
containment is vacuous.  Theorem~\ref{thm:main}'s hypothesis $q \ge 6$
is the joint requirement of:
\begin{itemize}[topsep=0.2em,itemsep=0.1em]
\item Theorem~\ref{thm:43-prefix} at $q' = q - 1 \ge 5$, i.e.,
$q \ge 6$ (for the $\mu_A$-recursion);
\item Theorem~\ref{thm:42-prefix} at $q' = q \ge 4$, i.e., $q \ge 4$
(for the $\Phi_h$-vanishing in Theorem~\ref{thm:reduction}).
\end{itemize}
The binding hypothesis is $q \ge 6$ (from the $\mu_A$-recursion).
The Phi-vanishing already kicks in at $q \ge 4$; equivalently the
reduction~\eqref{eq:reduction} holds at $q \ge 4$ but the prefix
bound only becomes positive at $q \ge 6$.
\end{remark}

\section{Sharpness}\label{sec:sharp}

\begin{lemma}[Sharpness reduction]\label{lem:sharp-reduction}
Suppose $B \subseteq E_{q-4}^-|_{V_\lambda}$ for $\lambda = (4, 4, 1^q)$
with $q \ge 6$.  Then every $T \in \supp(B)$ satisfies
$p(T) \ge q - 3 = n - 11$.
\end{lemma}

\begin{proof}
Set $j_0 := q - 4 = n - 12$.  Suppose for contradiction that
$T \in \supp(B)$ has $p(T) = n - 12$ exactly.  Then there exists
$w \in B$ with $C_T := [v_T]\,w \ne 0$, and by hypothesis
$T_{j_0}\,w = -w$.

\medskip\textbf{Locating entry $n - 11$ in $T$.}
By $p(T) = n - 12$, the cells $\{(1,1), \ldots, (n-12, 1)\}$ hold
entries $\{1, \ldots, n-12\}$, and $T(n-11, 1) \ne n - 11$.  Since
$q \ge 6$, $n - 12 = q - 4 \ge 2$, so row 1's column-1 cell $(1, 1)$
and row 2's column-1 cell $(2, 1)$ both sit in the prefix region.
The remaining $12$ cells of $\lambda$ are
\[
   \{(1,2),(1,3),(1,4),\;(2,2),(2,3),(2,4),\;(n-11,1),(n-10,1),(n-9,1),(n-8,1),(n-7,1),(n-6,1)\},
\]
holding the $12$ entries $\{n-11, n-10, \ldots, n\}$.

We claim entry $n - 11$ sits at $(1, 2)$:
\begin{itemize}[topsep=0.3em,itemsep=0.1em]
\item Not at $(n-10, 1), \ldots, (n-6, 1)$: column-$1$ strict
   increase from $T(n-12, 1) = n - 12$ forces $T(n-11, 1) > n - 12$,
   i.e., $T(n-11, 1) \in \{n-11, \ldots, n\}$.  For entry $n - 11$
   to sit at any $(k, 1)$ with $k \ge n-10$, column-$1$ monotonicity
   gives $T(n-11, 1) < T(k, 1) = n - 11$, but $T(n-11, 1) \ge n - 11$,
   contradiction.
\item Not at $(n-11, 1)$: this would mean $T(n-11, 1) = n - 11$,
   extending the prefix to $p(T) \ge n - 11 > n - 12$, contradicting
   $p(T) = n - 12$.
\item Not at $(1, 3)$: $T(1, 2) < T(1, 3) = n - 11$ would force
   $T(1, 2) < n - 11$, but all remaining entries are $\ge n - 11$,
   impossible.
\item Not at $(1, 4)$: $T(1, 3) < T(1, 4) = n - 11$ would force
   $T(1, 3) < n - 11$, impossible by the same argument.
\item Not at $(2, 2)$: $T(1, 2) < T(2, 2) = n - 11$ forces
   $T(1, 2) < n - 11$, impossible.
\item Not at $(2, 3)$: $T(2, 2) < T(2, 3) = n - 11$ forces
   $T(2, 2) < n - 11$, impossible.
\item Not at $(2, 4)$: $T(2, 3) < T(2, 4) = n - 11$ forces
   $T(2, 3) < n - 11$, impossible.
\end{itemize}
Hence $T(1, 2) = n - 11$.

\medskip\textbf{Off-diagonal swap.}
The cells of $j_0 = n - 12$ (at $(n - 12, 1)$) and $j_0 + 1 = n - 11$
(at $(1, 2)$) lie in different rows and columns.  Axial distance:
\[
   |d| \;=\; |c(1, 2) - c(n-12, 1)| \;=\; |1 - (1 - (n-12))| \;=\; n - 12 \ge 2
   \quad \text{for } q \ge 6.
\]
The Hoefsmit $2$-block at $T_{j_0}$ is non-degenerate: $T_{j_0}\,v_T
= \alpha\,v_T + \beta\,v_{T^*}$ with $\beta \ne 0$, where
$T^* := s_{j_0} T$ swaps entries $n-12$ and $n-11$ between cells
$(n-12, 1)$ and $(1, 2)$.

\medskip\textbf{$T^*$ is a valid SYT.}  Set $T^*(n-12, 1) = n - 11$
and $T^*(1, 2) = n - 12$, with all other entries identical to $T$.
Column-$1$ above: $T^*(n-13, 1) = n - 13 < n - 11 = T^*(n-12, 1)$
$\checkmark$ (when $n \ge 14$, i.e., $q \ge 6$).  Column-$1$ below:
$T^*(n-11, 1) = T(n-11, 1) \in \{n-10, \ldots, n\}$, so $T^*(n-12, 1)
= n - 11 < T^*(n-11, 1)$ $\checkmark$.  Row $1$: $T^*(1, 1) = 1 <
n - 12 = T^*(1, 2)$ $\checkmark$ (since $n - 12 \ge 2$ for $q \ge 6$);
$T^*(1, 2) = n - 12 < T(1, 3) \in \{n-10, \ldots, n\}$ $\checkmark$
(since $T(1, 3) \ne n-11$ by uniqueness of position of $n-11$ in $T$).
Column~$2$: $T^*(1, 2) = n - 12 < T(2, 2) \in \{n-10, \ldots, n\}$
$\checkmark$.  So $T^*$ is a valid SYT of $\lambda$.

\medskip\textbf{Prefix drop forces $T^* \notin \supp(B)$.}
$T^*(n-12, 1) = n - 11 \ne n - 12$, so $p(T^*) \le n - 13 < n - 12$.
By Theorem~\ref{thm:main}, every $U \in \supp(B)$ has $p(U) \ge
n - 12$, hence $T^* \notin \supp(B) \supseteq \supp(w)$.

\medskip\textbf{Coefficient contradiction.}  Compute
\[
   [v_{T^*}]\,(T_{j_0}\,w)
   \;=\; \sum_{U \in \supp(w)} C_U \cdot [v_{T^*}]\,(T_{j_0}\,v_U).
\]
The diagonal contribution $[v_{T^*}]\,T_{j_0}\,v_{T^*}$ requires
$T^* \in \supp(w)$ (no).  The off-diagonal contribution from $v_U$
requires $s_{j_0} U = T^*$, i.e., $U = s_{j_0} T^* = T$.  Hence
\[
   [v_{T^*}]\,T_{j_0}\,w \;=\; C_T \cdot \beta \;\ne\; 0,
\]
while $[v_{T^*}](-w) = -C_{T^*} = 0$.  This contradicts
$T_{j_0}\,w = -w$.  Hence no $T \in \supp(B)$ has $p(T) = n - 12$
exactly, and every such $T$ has $p(T) \ge n - 11$.
\qed
\end{proof}

\begin{theorem}[Sharpness, computationally verified]\label{thm:sharp}
For $\lambda = (4, 4, 1^q)$ at $q = 6$ (and conjecturally for $q \ge 6$),
$\supp(B)$ contains an SYT $T$ with $p(T) = n - 12 = q - 4$ exactly;
consequently $B \not\subseteq E_{q-4}^-|_{V_\lambda}$.
\end{theorem}

\begin{proof}
Existence at $q = 6$ is verified computationally
(\texttt{\textasciitilde/projects/scratch/2026-05-14-extended-sign-kill-44/explore\_44\_family.py}).
Combined with Lemma~\ref{lem:sharp-reduction}, $B \subseteq E_{q-4}^-$
would force every $U \in \supp(B)$ to have $p(U) \ge n - 11$,
contradicting the witness $T$ with $p(T) = n - 12$.  Hence
$B \not\subseteq E_{q-4}^-$.
\qed
\end{proof}

\section{Eigenspace summary}\label{sec:summary}

\begin{theorem}[Eigenspace summary, $\hatl = (4, 4)$]\label{thm:summary}
For $\lambda = (4, 4, 1^q)$ with $q \ge 6$,
\[
   B \;\subseteq\; E_\ell^+\;\cap\;\bigcap_{j=1}^{q-5} E_j^-|_{V_\lambda},
   \qquad
   B \not\subseteq E_{q-4}^-|_{V_\lambda}
   \quad (\text{at $q = 6$, conjectural at $q \ge 6$}).
\]
The $E_\ell^+$ containment is sharp with scalar action: $T_\ell|_B
= q\cdot \mathrm{Id}_B$ (May-13 evening Theorem A).
\end{theorem}

\begin{proof}
$E_j^-$ chain at $j \le q - 5$: Theorem~\ref{thm:main}.
Non-containment at $j = q - 4$ (at $q = 6$): Theorem~\ref{thm:sharp}.
$E_\ell^+$ containment with scalar action: May-13 evening (universal
across all $\lambda$ in stable regime).
\qed
\end{proof}

\section{Computational verification}\label{sec:verify}

\begin{theorem}[Verification]\label{thm:verify}
The structural results above were verified computationally at
$q \in \{1, 2, 3, 4, 5, 6\}$ (i.e., $n \in \{9, 10, 11, 12, 13, 14\}$)
mod $P = 100\,003$ at $Q = 11$.
\begin{itemize}[topsep=0.3em,itemsep=0.1em]
\item Two-piece decomposition $\dim \Rp_n V_\lambda = f^{\mu_A} +
   f^{\mu_h}$ verified at $q \in \{1, 2, 3\}$
   (Lemma~\ref{lem:Eplus-decomp}).
\item $\dim B = 5$ matching the May-13 multi-corner formula
   $f^{(3, 3)} = 5$ verified at $q \in \{4, 5, 6\}$ (stable regime).
   At $q = 2$ we get $\dim B = 10 = 2{\cdot}5$, and at $q = 3$ we get
   $\dim B = 15 = 3{\cdot}5$, both sub-stable.
\item $B \not\subseteq E_1^-$ at $q \in \{4, 5\}$
   (consistent with predicted threshold $\tau = q - 5 \le 0$):
   $(T_1+1)\,\Om v$ has $471, 523$ nonzero entries respectively for
   random $v$.
\item $B \subseteq E_1^-$ at $q = 6$ verified by rank probe of
   $(T_1+1)\,\Om$ on $5$ random vectors (all zero).
\item Sharpness witness at $q = 6$: an SYT $T \in \supp(B)$ with
   $p(T) = q - 4 = 2$ exactly:
   $T(1, \cdot) = (1, 3, 4, 5),\ T(2, \cdot) = (2, 6, 7, 8),\
   T(3..8, 1) = (9, 10, 11, 12, 13, 14)$.
\end{itemize}
\end{theorem}

\begin{proof}
Scripts in \texttt{\textasciitilde/projects/scratch/2026-05-14-extended-sign-kill-44/}:
\begin{itemize}[topsep=0.3em,itemsep=0.1em,leftmargin=2em]
\item \texttt{check\_decomp.py} --- two-piece $E_{n-1}^+$
   decomposition at $q \in \{1, 2, 3\}$.
\item \texttt{check\_dim\_B.py} --- $\dim B$ via full matrix
   computation at $q \in \{2, 3\}$ (showing sub-stable factor).
\item \texttt{verify\_E1minus.py} --- sparse-vector probe of
   $B \subseteq E_1^-$, dim probe, and sharpness witness search
   at $q \in \{4, 5, 6\}$.
\item \texttt{sparse\_ops.py} --- sparse Hecke operator application
   for vector-level Hoefsmit computation.
\end{itemize}
\qed
\end{proof}

\section{Discussion}\label{sec:discussion}

\subsection{What is new}

This paper closes the $\hatl = (4, 4)$ row of the extended sign-kill
conjecture at the predicted threshold $\tau = q - 5$, in the
containment direction (for $q \ge 6$), and reduces sharpness at
$j = q - 4$ to a boundary-prefix witness (verified at $q = 6$).
The recursion graph after this paper:
\begin{center}
\begin{tabular}{rl}
$(2, 2)$ & [base via support rule] \\
$\downarrow$ \\
$(3, 2),\;(4, 2)$ & [recurse on $(2, 2)$] \\
$\downarrow$ \\
$(3, 3)$ & [recurse on $(3, 2)$ for $\mu_A$, hook for $\mu_h$] \\
$\downarrow$ \\
$(4, 3)$ & [recurse on $(3, 3)$, $(4, 2)$ for contributing branches;
            $(3, 2)$, hook for vanishing] \\
$\downarrow$ \\
$\mathbf{(4, 4)}$ & [recurse on $(4, 3)$ for $\mu_A$; $(4, 2)$ for $\mu_h$ vanishing]
\end{tabular}
\end{center}

\subsection{$\hatl = (4, 4)$ vs $\hatl = (3, 3)$}

Both shapes are $r = 1$, so both have a two-piece $E_{n-1}^+$
decomposition.  The structural similarity is that the same-row
piece is the one that vanishes, while the doubly-stripped piece
carries the dimension.

The substantive difference:
\begin{itemize}[topsep=0.2em,itemsep=0.1em]
\item For $(3, 3)$, the same-row inner shape $\mu_h = (3, 1^{q+1})$
is a \emph{hook}; vanishing comes from the May-12 hook column-$1$
theorem.
\item For $(4, 4)$, the same-row inner shape $\mu_h = (4, 2, 1^q)$
is \emph{not} a hook (it has a length-$2$ second row); vanishing
requires the $(4, 2)$ extended sign-kill theorem.
\end{itemize}

This is the first case in the $r = 1$ stratum where the same-row
vanishing requires a non-hook input.  The general principle ---
$\Rp^{(\mu_h)}_{\ell+1}\,(\text{appropriate inner $B$}) \subseteq E_1^-$,
killed by the extra rightmost $\Rp^{(\mu_h)}_{n-6}$ factor --- carries
over uniformly: only the specific extended sign-kill theorem to invoke
shifts.

\subsection{The recursion ladder reaches $|\hatl| = 8$}

The proved table after this paper:
\begin{center}
\begin{tabular}{c|c|c}
$\hatl$ & predicted $\tau$ & status \\
\hline
$(2, 2)$ & $q - 1$ & proved (May-14 fourth pass) \\
$(3, 2)$ & $q - 2$ & proved (May-14 fifth pass) \\
$(3, 3)$ & $q - 3$ & proved (May-14 seventh pass) \\
$(4, 2)$ & $q - 3$ & proved (May-14 sixth pass; same-row gap to close) \\
$(4, 3)$ & $q - 4$ & proved (May-14 eighth pass) \\
$(4, 4)$ & $q - 5$ & proved (this paper) \\
\hline
$(5, 2)$ & $q - 4$ & open (recurses on $(4, 2)$ + hook) \\
$(5, 3)$ & $q - 5$ & open (three corner pairs; needs $(4, 2), (5, 2), (4, 3)$) \\
\end{tabular}
\end{center}
The rows $|\hatl| \le 8$ are now essentially closed.  The natural
next target is $\hatl = (5, 3)$ at $|\hatl| = 8$, which has three
corner pairs (so a four-piece $E_{n-1}^+$ decomposition with three
non-hook branches), and would need the $(5, 2)$ extended sign-kill
theorem (currently open) as an input.

\subsection{Looking ahead}

A general structural pattern emerging: in the $r = 1$ stratum the
reduction has two pieces (doubly-stripped + same-row at $c_1$);
in the $r = 2$ stratum the reduction has four pieces (three corner
pairs + same-row at $c_1$).  Whether a branch contributes vs vanishes
depends on the inner shape's $\ell + 1$ value relative to $\lambda$'s
chain length.  The contributors' inputs are extended sign-kill
theorems at smaller $\hatl$; the vanishers' inputs are either hook
column-$1$ or $(4, 2)$-type extended sign-kill, depending on whether
the same-row inner shape is a hook.

\subsection{What this is not}

This paper continues the structural sharpening of $B$ via its
$T_j$-eigenspace containments --- the programme orthogonal to (and
unaffected by) the iso-programme refutations of May-13 and May-14.
We learn more about \emph{where} $B$ sits inside $V_\lambda$ via
these eigenspace conditions, without claiming $B$ equals any standard
representation-theoretic object.

\section{Status}\label{sec:status}

\paragraph{Established (this paper).}
\begin{itemize}[topsep=0.3em,itemsep=0.1em]
\item Lemma~\ref{lem:Eplus-decomp}: two-piece decomposition of
   $E_{n-1}^+$.
\item Theorem~\ref{thm:reduction}: single-branch reduction (after the
   same-row branch vanishes), $\lambda = (4, 4, 1^q)$, $q \ge 4$.
\item Theorem~\ref{thm:main}: $B \subseteq E_j^-$ for $j \le q - 5$
   at $q \ge 6$.
\item Lemma~\ref{lem:sharp-reduction}: structural sharpness
   reduction at $j = q - 4$ for $q \ge 6$.
\item Theorem~\ref{thm:sharp}: sharpness at $j = q - 4$ for $q = 6$
   (computational witness).
\end{itemize}

\paragraph{Open.}
\begin{itemize}[topsep=0.3em,itemsep=0.1em]
\item Structural (non-computational) sharpness witness across all
   $q \ge 6$.
\item The $(5, 3)$ family (and beyond) at $|\hatl| = 8$.
\item The $(4, 2)$ same-row gap (Remark 3.5 of seventh-pass paper):
   pending companion erratum.
\end{itemize}

\paragraph{Push status.}
Pending PAT refresh; unpushed-commit count on
\texttt{clio-vega/proofs}: 70 (after this paper).

\end{document}
